\section{Verification}
\label{sec:verification}

The C++ in \texttt{core/} and the equations in this document must stay in agreement. We pin them together with two layers of tests: doctest unit tests in \texttt{core/tests/} that run on every CI build, and the Jupyter notebook \texttt{analysis/dynamics\_tour.ipynb} that exercises the same code from Python with plots and looser numerical tolerances. This section maps each result derived above to the test that verifies it.

\subsection{Test crosswalk}

\begin{center}
\begin{tabular}{l l l}
\toprule
\textbf{Notebook test} & \textbf{Verifies} & \textbf{Equation} \\
\midrule
1. Motor-lag step response  & First-order lag on $\dot v$                                & \eqref{eq:vdot} \\
2. Pure rotation            & No side slip ($v = 0$, $p_x = p_y = 0$)                    & \eqref{eq:noslip} \\
                            & Yaw rate from differential commands                        & \eqref{eq:fwd-kin} \\
3. Constant arc             & Steady-state turn radius $R = v_{ss}/\omega_{ss}$          & \eqref{eq:icr} \\
4. Linearisation residual   & Analytic $A$ matches the nonlinear flow to first order    & \eqref{eq:Amat}, \eqref{eq:lin-residual} \\
\bottomrule
\end{tabular}
\end{center}

The \texttt{core/tests/} doctest binary runs three lower-level checks: state derivative at the origin under symmetric control, single-step RK4 sanity, and finiteness of $A$ and $B$. It is the build-time canary; the notebook is the analyst-facing companion that produces plots.

\subsection{Tests still to add}

Two checks are derived in this document but not yet wired into either test layer:

\begin{enumerate}
    \item \textbf{Constraint satisfaction along long rollouts.} \eqref{eq:noslip} is identically satisfied by the parameterisation \eqref{eq:state-ode}, so this is essentially a regression test against the position equations being edited. Future Test~5 of the notebook: roll out for several seconds under a varying control sequence and assert
        \[
            \bigl|-\sin\theta(t)\,\dot p_x(t) + \cos\theta(t)\,\dot p_y(t)\bigr| < 10^{-12}
        \]
        at every time step.
    \item \textbf{Finite-difference Jacobian comparison in C++.} \S\ref{sec:linearization} derives $A$ and $B$ analytically; the doctest binary should compare them to centred-difference estimates at randomly sampled operating points. This is the cheapest of the three verification routes from \S\ref{sec:linearization} and the highest-value to add first.
\end{enumerate}

When either of these is added, this section's table grows by one row and the corresponding bullet here is removed.

\subsection{Editing discipline}

The chain of authority across the project is:
\begin{quote}
    LaTeX (\textit{here}) $\;\longrightarrow\;$ C++ (\texttt{core/}) $\;\longrightarrow\;$ Python (\texttt{bindings/}, notebook).
\end{quote}

If the LaTeX changes, the C++ must change to match before the doctest binary is re-run. If the C++ changes without the LaTeX, the doctest will likely still pass --- but the document will be wrong, and the next person to read it will derive the wrong control law from the wrong equation. The reviewer's job on any \texttt{core/} pull request is to confirm the LaTeX is updated in lockstep.
